% ===========================================================
%  第 3 章 矩阵运算 —— 高等代数【深化】拓展集
%  来源：A. 北大《高等代数》第五版 第四章 矩阵 §3–§7
%          （PDF p129–p144 = 印刷页 117–132，偏移 PDF = 印刷页 + 12，已实测页脚）
%        B. 樊启斌《高等代数典型问题与方法》§4.1、§4.3
%          （PDF p202–p203 = 印刷页 193–194；PDF p223–p228 = 印刷页 214–219，
%            偏移 PDF = 印刷页 + 9）
%  筛选依据：AGENTS.md §2 决定二「只补能加深理解、反哺解题的深化知识」，
%            服务考纲「二、矩阵」：线性运算、乘法、方阵的幂、乘积的行列式、
%            逆矩阵、伴随矩阵、初等变换、初等矩阵、矩阵的等价、分块矩阵。
% ===========================================================
%
%  【保留 17 块 / 四组】
%   深化一 乘法与逆的“非数”特性（对应深化清单 3.1）：
%     ① 乘法不可交换是常态（樊 4.1.3 印刷 193 + 樊例 4.1 印刷 193）
%     ② 消去律、幂公式都与数的运算不同（樊 4.1.1/4.1.2 印刷 193）
%     ③ 方阵乘积的行列式与可逆性（北大 §3 印刷 117 + 樊 4.1.2 印刷 193）
%   深化二 逆矩阵的等价刻画网（对应深化清单 3.2）：
%     ④ 可逆的等价刻画网（北大 §6 印刷 128 + §4 印刷 119；樊 4.3.1 印刷 214）
%     ⑤ 只验一边即可判定互逆（北大 §4 印刷 120，系第 2 章同一段的不同用法）
%   深化三 初等矩阵 = 初等变换的坐标表示（对应深化清单 3.3）：
%     ⑥ 初等矩阵的定义与三类形状（北大 §6 印刷 125 + 樊 4.3.1 印刷 214）
%     ⑦ 左乘做行变换、右乘做列变换（北大 §6 印刷 126；樊 4.3.1 印刷 214）
%     ⑧ 等价标准形与“1 的个数 = 秩”（北大 §6 定理 5 印刷 127；樊 4.3.1 印刷 214）
%     ⑨ A 可逆 ⟺ A 是初等矩阵的乘积（北大 §6 定理 6 印刷 128）
%     ⑩ 等价判据 r(A)=r(B) 与 (A|E) 求逆的原理（北大 §6/§7 印刷 128、129）
%   深化四 分块矩阵与秩的乘法不等式（对应深化清单 3.4、3.5）：
%     ⑪ 分块矩阵的运算律（北大 §5 印刷 122；樊 4.3.2 印刷 214）
%     ⑫ 准对角矩阵的行列式与逆（北大 §5 印刷 124/125；樊 4.3.2 印刷 215）
%     ⑬ 分块上三角求逆与 Schur 补（北大 §5 印刷 123、124）
%     ⑭ 分块初等变换：左乘做分块行变换（北大 §7 印刷 130）
%     ⑮ 分块消去法：一行一列地造零块（北大 §7 例 3、例 4 印刷 131、132）
%     ⑯ 秩的乘法不等式 r(AB)≤min（北大 §3 定理 2 印刷 117、118）
%     ⑰ 可逆矩阵不改变秩（北大 §3 定理 4 印刷 121）
%
%  【跳过及理由】
%   · beida p129/p130 顶部「例 4」「例 3」、p133 的分块乘法算例、p137 的准对角算例、
%     p140 例 1、p141 例 2、p142 例 1/例 2、p143 例 3 → 全部为例题/计算演示，按深化原则跳过。
%   · beida p144「习 题」1.–… → 习题，跳过。（p144 上半页保留：LU 型归约的证法）
%   · beida p135「D=[[A,O],[C,B]] 的逆」→ 与第 2 章已收录的同名结论重复，且不及
%     p123–124 的推导完整，取后者。
%   · beida p131「定义 8 逆矩阵」「定义 9 伴随矩阵」→ 已由第 2 章 deep_ch02_beida.tex
%     收录（逆矩阵定义与唯一性、伴随矩阵的定义），本章不重复。
%   · beida p132「逆矩阵的运算性质（转置、乘积）」「用逆矩阵推导克拉默法则」
%     → 已由第 2 章收录，本章不重复。
%   · 樊 4.2 矩阵的伴随、4.4 矩阵的秩的其余版面 → 不属本次指定页范围。
%   · 樊 p202 例 4.2、p203 例 4.3/例 4.4、p224–p228 例 4.43–4.51
%     → 例题，按深化原则跳过（仅从“注”中提取可复用的公式/方法，见 ⑩ 与 ③）。
%
%  【体例说明】
%   · 每块首行说明来源页（北大 §x = 北大第四章 §x；樊 x.y = 樊书 §x.y）。
%   · 每块正文末以「【反哺点】…」收束，写明它服务考纲「二、矩阵」的哪一条要求、
%     以及可迁移到哪类题；其中 4 块（④ 可逆等价网、⑧ 等价标准形、⑭ 分块初等变换、
%     ⑯ 秩的乘法不等式）的【反哺点】另起一段、其后还有 2–4 行具体用法，属该段的展开。
%   · 每块正文 557–1060 字符（多数 750–900），符合"每块聚焦一个结论、讲清为什么"的定位。
% ===========================================================


\fsec{深化一　矩阵乘法与逆的「非数」特性}

% ---------------- ① 不可交换 ----------------
\begin{fdeep}{矩阵乘法不可交换是常态，「与 $A$ 可交换」是有结构的一族}
矩阵乘法满足结合律 $(AB)C=A(BC)$、左（右）分配律 $A(B+C)=AB+AC$，但\textbf{一般不满足交换律}：
$AB$ 与 $BA$ 未必相等，甚至可能一个是 $m\times m$、另一个是 $n\times n$，连阶数都不同（例如
$A$ 为 $m\times n$、$B$ 为 $n\times m$）。全部 $n$ 阶方阵中满足 $AB=BA$ 的 $(A,B)$ 只是一小部分。

可交换矩阵的三条结构性质（樊 4.1.3）：
\begin{enumerate}[leftmargin=2.0em,itemsep=0.2em]
\item 若 $AB=BA$，则称 $A$ 与 $B$ \textbf{可交换}，可交换是矩阵之间的一种关系；
\item \textbf{一个 $n$ 阶方阵是数量矩阵（$\lambda E$）当且仅当它与任意 $n$ 阶方阵都可交换}；
\item $A=\operatorname{diag}(d_1,d_2,\dots,d_n)$，当 $i\ne j$ 时 $d_i\ne d_j$，则与 $A$ 可交换的矩阵只能是（同阶）\textbf{对角矩阵}。
\end{enumerate}
第 ③ 条的来法只有一行：设 $B=(b_{ij})$，比较 $AB$ 与 $BA$ 的 $(i,j)$ 元得
$(d_i-d_j)b_{ij}=0$，于是 $i\ne j$ 时 $b_{ij}=0$。

考纲只说"运算规律"，而上表说明\textbf{哪些规律没了}：交换律没有。由此立刻得到两条考场判据——
（a）凡是要用 $AB=BA$ 的推导都必须先说明理由（数量矩阵、对角矩阵、多项式于同一矩阵等）；
（b）"求与 $A$ 可交换的全部矩阵"这一类题有统一解法：\textbf{设未知矩阵、写 $AB=BA$、逐元素比对}
（对角可交换题更是一行出结果）。注意第 ③ 条要求 $d_i$ 互异：若有重根，解集是相应尺寸的块对角矩阵。
【反哺点】服务考纲「矩阵」中"掌握矩阵的线性运算、乘法、转置以及它们的运算规律"。
\end{fdeep}

% ---------------- ② 消去律与幂公式 ----------------
\begin{fdeep}{消去律与幂公式：哪些能照搬，哪些会出事}
（樊 4.1.1、4.1.2）设 $A,B,C$ 为矩阵，$\lambda,\mu$ 为数。
\textbf{照搬成立的}：
\fml{\lambda(AB)=(\lambda A)B=A(\lambda B),\qquad A(B+C)=AB+AC,\qquad (AB)^{\mathrm{T}}=B^{\mathrm{T}}A^{\mathrm{T}}}
\fml{A^{m}A^{n}=A^{m+n},\qquad (A^{m})^{n}=A^{mn},\qquad \operatorname{tr}(AB)=\operatorname{tr}(BA)}
\textbf{不能照搬的}：由 $AB=AC$ 且 $A\ne O$ \textbf{推不出} $B=C$。
只有当 $A$ 可逆（或 $A$ 为可逆因子链中一员）时才可约去：由 $AB=AC$ 且 $A$ 可逆得 $B=C$。

由此派生的一类典型错误（樊例 4.1）：数的立方差公式 $A^3-B^3=(A-B)(A^2+AB+B^2)$ 对矩阵\textbf{不真}。
因为展开后
\fml{(A-B)(A^2+AB+B^2)=A^3+A^2B+AB^2-BA^2-BAB-B^3}
中间那些项只有在 $AB=BA$ 时才相消；一般情形必须补成
\fml{A^3-B^3=(A-B)(A^2+AB+B^2)\ \Longleftrightarrow\ A^2B+AB^2=BA^2+BAB,}
而这个条件一般\emph{不成立}。取 $A=\begin{pmatrix}1&0\\0&0\end{pmatrix},
B=\begin{pmatrix}0&1\\0&0\end{pmatrix}$ 即得 $AB\ne BA$ 的反例。

凡是把数的乘法公式（平方差、立方差、二项式展开）搬到矩阵上的题，
\textbf{判定可否使用的唯一标准就是各因子是否两两可交换}；两两可交换时二项式定理也对，
否则必须老老实实逐项展开。这既是选择题设陷阱的方式，也是"判断某等式是否恒成立"的答题模板。
【反哺点】服务考纲「矩阵」中"矩阵的线性运算、乘法、转置以及它们的运算规律、方阵的幂"。
\end{fdeep}

% ---------------- ③ 乘积的行列式与可逆性 ----------------
\begin{fdeep}{方阵乘积的行列式与可逆性：退化性「可遗传、可传染」}
\textbf{定理 1}（北大 §3）：设 $A,B$ 是数域 $P$ 上两个 $n\times n$ 矩阵，则
\fml{|AB|=|A|\,|B|.}
用数学归纳法不难推广到多个因子的情形：$|A_1A_2\cdots A_m|=|A_1||A_2|\cdots|A_m|$。

\textbf{定义 6}（北大 §3）：数域 $P$ 上的 $n\times n$ 矩阵 $A$ 称为\textbf{非退化}的，如果 $|A|\ne 0$；
否则称为\textbf{退化}的。显然，一个 $n\times n$ 矩阵是非退化的充分必要条件为它的秩等于 $n$。

\textbf{推论 1}（北大 §3）：设 $A,B$ 是数域 $P$ 上 $n\times n$ 矩阵，矩阵 $AB$ \textbf{为退化}的
充分必要条件是 $A,B$ 中至少有一个是退化的。

证明只用到乘法定理取反：$|AB|=|A||B|$，而 $|AB|=0\iff |A|=0$ 或 $|B|=0$。

（a）\textbf{方阵乘积 $AB$ 可逆 $\iff$ $A$ 与 $B$ 都可逆}，且此时 $(AB)^{-1}=B^{-1}A^{-1}$——
凡是"某乘积可逆 / 不可逆"的选择题，直接拆成因子讨论，不必求逆；
（b）\textbf{$|AB|=|A||B|$ 可用于"算不出来就换序"：}$|AB|=|BA|$ 对同阶方阵无条件成立，
这是抽象行列式（给 $AB$ 的阶数远小于 $A,B$ 阶数）最常用的降阶通道。
（樊 4.1 的"注"另给两个提速公式，小阶计算可直接用：
$\begin{pmatrix}a&b\\c&d\end{pmatrix}^{*}=\begin{pmatrix}d&-b\\-c&a\end{pmatrix}$，
$\begin{pmatrix}A&O\\O&B\end{pmatrix}^{-1}=\begin{pmatrix}A^{-1}&O\\O&B^{-1}\end{pmatrix}$。）
【反哺点】服务考纲「矩阵」中"掌握方阵乘积的行列式"。由此立刻得到两条考场上极常用的判据：
\end{fdeep}


\fsec{深化二　逆矩阵的等价刻画网}

% ---------------- ④ 等价刻画网 ----------------
\begin{fdeep}{可逆的等价刻画网：一张图串起整个线性代数}
对 $n$ 阶方阵 $A$，以下命题\textbf{两两等价}：
\begin{enumerate}[leftmargin=2.0em,itemsep=0.18em]
\item $A$ 可逆（存在 $B$ 使 $AB=BA=E$）；
\item $|A|\ne 0$，即 $A$ \textbf{非退化}（北大 §4 定理 3）；
\item $\operatorname{r}(A)=n$（行秩 $=$ 列秩 $=n$）；
\item $A$ 的 $n$ 个行向量线性无关（等价地：$n$ 个列向量线性无关）；
\item 齐次方程组 $AX=0$ \textbf{只有零解}（即解空间维数为 $0$）；
\item $A$ 与单位矩阵 $E$ \textbf{等价}，即存在初等矩阵 $P_1,\dots,P_k$、$Q_1,\dots,Q_l$ 使
      $A=P_1\cdots P_kE\,Q_1\cdots Q_l$（此时可取 $l=0$，见 ⑨）；等价地 $A$ 的等价标准形为 $E$；
\item $A$ 的特征值全不为 $0$（第 5 讲引入特征值后补入本网）；
\item $A$ 的伴随矩阵 $A^{*}$ 可逆（因 $|A^{*}|=|A|^{\,n-1}$，$n\ge2$）。
\end{enumerate}
\textbf{这张网为什么能连起来}：②$\Rightarrow$①由 $A\bigl(\tfrac{1}{|A|}A^{*}\bigr)=E$ 得到，①$\Rightarrow$②
由 $|A||A^{-1}|=|E|=1$ 得到；①$\Rightarrow$⑥由 $\operatorname{r}(A)=n$ 与定理 5（见 ⑧）得到，
⑥$\Rightarrow$①由定理 6（见 ⑨）得到。所以整张网只需\textbf{两条独立的事实}
（行列式判据、初等变换保秩）就能串通。
【反哺点】服务考纲「矩阵」中"理解矩阵可逆的充分必要条件"——这是考纲原话，也是整门课的中心枢纽。
考试中的实际价值在于\textbf{换一条最好走的路}：要证可逆，优先找 $|A|\ne0$ 或 $\operatorname{r}(A)=n$；
要给方程组"只有零解"，就看 $|A|\ne0$；要说明某行列式非零，可以改证"该矩阵的列线性无关"。
题目把条件包装成哪一种说法，就从网的哪一格切进去。
\end{fdeep}

% ---------------- ⑤ 只验一边 ----------------
\begin{fdeep}{只验一边即可判定互逆（$AB=E$ 的信息量）}
北大 §4 定理 3 之后的一句断言：对于 $n$ 阶方阵 $A,B$，如果
\fml{AB=E,}
那么 $A,B$ 就都是可逆的，并且它们互为逆矩阵。

这不同于一般代数系统的要求——定义里写的是 $AB=BA=E$ 两个等式；但对\textbf{同阶方阵}，
一个等式就足够了（依据是 $|AB|=|A||B|$：由 $|A||B|=1$ 得 $|A|\ne0$，再用 $A^{-1}$ 左乘）。
与之配套的是定理 3 给出的逆矩阵公式与行列式关系：
\fml{A^{-1}=\frac{1}{d}A^{*}\quad(d=|A|\ne0),\qquad |A^{-1}|=d^{-1}.}

（a）判定互逆或证明"$B$ 就是 $A^{-1}$"时，\textbf{只需验证 $AB=E$}，不必再验 $BA=E$——
在证明题里常常能省掉一半书写（这也解释了为什么矩阵方程 $AX=B$ 解出 $X=A^{-1}B$ 后可以不作回代检验）；
（b）注意该结论对\textbf{方阵}才成立：$A$ 为 $m\times n$、$B$ 为 $n\times m$（$m\ne n$）时
$AB=E_m$ 与 $BA=E_n$ 不能互相推出，这是"单侧逆"的常见陷阱。
【反哺点】服务考纲「矩阵」中"掌握逆矩阵的性质"。两条实战用法：
\end{fdeep}


\fsec{深化三　初等矩阵 = 初等变换的坐标表示}

% ---------------- ⑥ 定义与三类形状 ----------------
\begin{fdeep}{初等矩阵的定义与三类形状（定义 10）}
（北大 §6）对单位矩阵 $E$ 作\textbf{一次}初等变换所得到的矩阵称为\textbf{初等矩阵}。
显然，初等矩阵都是方阵，每个初等变换都有一个与之相应的初等矩阵，共三类，分别记为
$P(i,j)$（把 $E$ 的第 $i$ 行与第 $j$ 行互换：除 $(i,j)$ 与 $(j,i)$ 两位为 $1$、$(i,i)$ 与 $(j,j)$ 两位为 $0$ 外，
其余与 $E$ 相同）、$P(i(c))$（用数域 $P$ 中非零数 $c$ 乘 $E$ 的第 $i$ 行，即对角元 $1$ 换成 $c$）、
$P(i,j(k))$（把 $E$ 的第 $j$ 行的 $k$ 倍加到第 $i$ 行上，即 $(i,j)$ 位置放 $k$）。
三类矩阵就是全部的初等矩阵：对单位矩阵作一次初等\textbf{列}变换得到的也在这三类之中，
例如把 $E$ 的第 $i$ 列的 $k$ 倍加到第 $j$ 列，仍得到 $P(i,j(k))$。

\textbf{一一对应}记忆，并牢记第三条：\textbf{用 $c$ 乘某行要 $c\ne0$}（否则不可逆），
而"把某行的 $k$ 倍加到另一行"的 $k$ \textbf{可以为零}（$\Rightarrow$ $P(i,j(k))$ 允许 $k=0$，此时它是 $E$）。
凡要求把某矩阵分解成初等矩阵乘积的题，第一步永远是"把行变换过程逐条翻译成 $P$ 的乘积"。
【反哺点】服务考纲「矩阵」中"理解初等矩阵的概念"。关键在于把三类初等矩阵与三类初等变换
\end{fdeep}

% ---------------- ⑦ 左行右列 ----------------
\begin{fdeep}{左乘做行变换、右乘做列变换（§6 引理）}
\textbf{引理}（北大 §6）：对一个 $s\times n$ 矩阵 $A$ 作一次初等行变换，就相当于在 $A$ 的\textbf{左}边
乘上相应的 $s\times s$ 初等矩阵；对 $A$ 作一次初等列变换，就相当于在 $A$ 的\textbf{右}边乘上相应的
$n\times n$ 初等矩阵。

证明思路（只写行变换）：设 $B=(b_{ij})$ 为任意 $s\times s$ 矩阵，$A_1,A_2,\dots,A_s$ 为 $A$ 的行向量，
由分块乘法得 $BA$ 的第 $i$ 行为 $b_{i1}A_1+b_{i2}A_2+\cdots+b_{is}A_s$。
特别地取 $B=P(i,j)$，则 $BA$ 的第 $i$ 行是 $A_j$、第 $j$ 行是 $A_i$——交换了两行；
取 $B=P(i(c))$，得 $BA$ 的第 $i$ 行为 $cA_i$；取 $B=P(i,j(k))$，得第 $i$ 行为 $A_i+kA_j$。
列变换的情形转置后同理（此时 $A$ 的列向量被 $B$ 的列线性组合）。

这条引理是全部初等变换算法的\textbf{合法性来源}：它把"对矩阵做操作"改写成"乘一个矩阵"，
于是"$A$ 与 $B$ 等价"就变成"存在可逆 $P,Q$ 使 $PAQ=B$"。
记忆口诀："\textbf{左行右列}——$P$ 左乘动行，$Q$ 右乘动列"。它还能解释两个易错点：
（a）求逆时 $(A\mid E)\to(E\mid A^{-1})$ 只能做\textbf{行}变换（因为全程是 $P_k\cdots P_1A=E$）；
（b）若对该分块做列变换，就破坏了右边的 $E$，得到的不是 $A^{-1}$。
【反哺点】服务考纲「矩阵」中"理解初等矩阵"与"掌握用初等变换求矩阵的秩和逆矩阵的方法"。
\end{fdeep}

% ---------------- ⑧ 等价标准形 ----------------
\begin{fdeep}{等价标准形与「主对角线上 1 的个数 = 秩」（定理 5）}
\textbf{定义 11}（北大 §6）：矩阵 $A$ 与 $B$ 称为\textbf{等价}的，如果 $B$ 可以由 $A$ 经过一系列初等变换得到。
等价是矩阵间的一种关系，不难证明它具有\textbf{自反性、对称性和传递性}。

\textbf{定理 5}：任意一个 $s\times n$ 矩阵 $A$ 都与一形式为
\fml{\begin{pmatrix}1&&&&&\\ &1&&&&\\ &&\ddots&&&\\ &&&1&&\\ &&&&\ddots&\\ &&&&&0\end{pmatrix}}
的矩阵等价，它称为 $A$ 的\textbf{标准形}，主对角线上 $1$ 的个数等于 $A$ 的秩（$1$ 的个数可以是零）。

证明思路（消去法，就是高斯消元的矩阵版）：$A=O$ 时已是标准形，故设 $A\ne O$。
当 $a_{11}\ne0$ 时，把其余行减去第一行的 $a_{11}^{-1}a_{i1}$（$i=2,\dots,s$）倍，
再把其余列减去第一列的 $a_{11}^{-1}a_{1j}$（$j=2,\dots,n$）倍，最后用 $a_{11}^{-1}$ 乘第一行，
$A$ 就变成 $\begin{pmatrix}1&0\\0&A_1\end{pmatrix}$（$A_1$ 是 $(s-1)\times(n-1)$ 矩阵）；
对 $A_1$ 重复同样步骤，归纳即得标准形。末句依据：初等变换不改变矩阵的秩，
而标准形的秩就是对角线上 $1$ 的个数。
【反哺点】服务考纲「矩阵」中"了解矩阵等价的概念"，并直接服务于"掌握用初等变换求矩阵的秩"。
这段话给出两个考点：\textbf{（a）标准形是等价类的唯一标签}——标准形由秩唯一决定，
所以 $A$ 与 $B$ 等价 $\iff$ 它们的标准形相同 $\iff$ $\operatorname{r}(A)=\operatorname{r}(B)$；
\textbf{（b）求秩的合法性}——正因为初等变换不改变秩，才可以用行变换把矩阵化成阶梯形数非零行。
另外要严格区分：等价（相抵）只需左右各乘可逆矩阵，\textbf{相似要求 $Q=P^{-1}$、合同要求 $Q=P^{\mathrm{T}}$}，
三者不能混用（相似、合同属第 5、6 节）。
\end{fdeep}

% ---------------- ⑨ 可逆 = 初等矩阵之积 ----------------
\begin{fdeep}{$A$ 可逆 $\iff$ $A$ 可表成初等矩阵的乘积（定理 6）}
把 ⑧ 的结论写细：$A,B$ 等价的充分必要条件是存在可逆的 $s$ 阶矩阵 $P$ 与可逆的 $n$ 阶矩阵 $Q$ 使
\fml{A=PBQ.}
因为对 $A$ 作初等变换就相当于用相应的初等矩阵去乘 $A$，于是 $A,B$ 等价的充要条件还可写成：
存在初等矩阵 $P_1,\dots,P_k,Q_1,\dots,Q_l$ 使
\fml{A=P_1P_2\cdots P_k\,B\,Q_1Q_2\cdots Q_l.}

\textbf{定理 6}（北大 §6）：$n$ 阶矩阵 $A$ 为可逆的充分必要条件是它能表成一些初等矩阵的乘积：
\fml{A=Q_1Q_2\cdots Q_m.}
\textbf{推论 2}：可逆矩阵总可以经过一系列初等行变换化成单位矩阵。

证明思路：$A$ 的秩为 $n$，由上一条定理它的标准形是 $E$，故有初等矩阵
$P_1,\dots,P_m,Q_1,\dots,Q_r$ 使 $P_1\cdots P_m\,A\,Q_1\cdots Q_r=E$；
把 $Q$ 全部移到右边（$A$ 可逆 $\Rightarrow$ 右端是初等矩阵之积，仍是可逆矩阵），即得结论。
推论 2 是 $l=0$ 的特殊情形（只用行变换），也正是 $(A\mid E)\to(E\mid A^{-1})$ 的合法性依据。

这条定理是把"可逆"这个\textbf{纯代数条件}翻译成"初等变换可消成 $E$"这个\textbf{可操作条件}的桥梁，
从而把求逆从"算 $n^2$ 个代数余子式"变成"做行变换"，计算量从 $O(n!)$ 级降到多项式级。
考场上遇到"证明某矩阵可表成初等矩阵乘积"或"讨论初等矩阵之积的性质"，都回到这条定理；
而推论 2 则说明：\textbf{可逆矩阵一定可以只用行变换化成 $E$}——这是 $(A\mid E)$ 法的理论保证。
【反哺点】服务考纲「矩阵」中"理解初等矩阵"与"掌握用初等变换求逆矩阵的方法"。
\end{fdeep}

% ---------------- ⑩ 等价判据与 (A|E) 法的原理 ----------------
\begin{fdeep}{等价判据 $\operatorname{r}(A)=\operatorname{r}(B)$ 与 $(A\mid E)$ 求逆的原理（(6) 式）}
樊 4.3.1 的要点：（1）$A$ 可逆 $\iff$ 存在有限个初等矩阵 $P_1,P_2,\dots,P_l$ 使 $A=P_1P_2\cdots P_l$；
（2）$A,B$ 均为 $m\times n$ 矩阵，则 $A$ 与 $B$ 等价 $\iff$ 存在 $m$ 阶可逆矩阵 $P$ 与 $n$ 阶可逆矩阵 $Q$ 使
\fml{PAQ=B;}
（3）$m\times n$ 矩阵 $A$ 的秩为 $r$ 的充要条件是存在 $m$ 阶可逆矩阵 $P$ 与 $n$ 阶可逆矩阵 $Q$ 使
\fml{PAQ=\begin{pmatrix}E_r&O\\ O&O\end{pmatrix}\qquad(E_r\ \text{为}\ r\ \text{阶单位矩阵}),}
即\textbf{等价标准形把"秩"翻译成"可逆矩阵能消掉多少"：}一个矩阵的秩，就是它所能化到的标准形中 $1$ 的个数。

北大 §6 末把求逆过程写成矩阵形式：把 $A$ 与 $E$ 凑成 $n\times2n$ 矩阵 $(A\ E)$，按分块乘法，求逆过程即
\fml{P_m\cdots P_1\,(A\ E)=(P_m\cdots P_1A\ \ P_m\cdots P_1E)=(E\ \ A^{-1}).\qquad(6)}

（a）等价判据"秩相等"把"存在可逆 $P,Q$"这一抽象条件化为\textbf{一个数相等}，选择题可直接秒判；
（b）(6) 式把 $(A\mid E)\xrightarrow{\text{行变换}}(E\mid A^{-1})$ 解释清楚了：
化成 $E$ 的\textbf{同一个行变换序列}同时作用在右边的 $E$ 上，就得到 $A^{-1}$——
所以\textbf{只能在右边拼接 $E$、只能做行变换}，不能拼在上方、也不能掺列变换。
【反哺点】服务考纲「矩阵」中"了解矩阵等价"与"掌握用初等变换求矩阵的秩和逆矩阵的方法"。
\end{fdeep}


\fsec{深化四　分块矩阵与秩的乘法不等式}

% ---------------- ⑪ 运算律 ----------------
\begin{fdeep}{分块矩阵的运算律：什么时候能把「块」当「数」}
（北大 §5、樊 4.3.2）把 $m\times n$ 矩阵用横线与竖线分成若干"子块"，从而将原矩阵看作由若干小矩阵组成，
即所谓\textbf{分块矩阵}。分块之后可以简化运算，也便于看出矩阵间的关系。设
\fml{A=(a_{ij})_{m\times n}\ \text{分块为}\ A=(A_{pq})_{r\times t},\qquad B=(b_{ij})_{m\times n}\ \text{分块为}\ B=(B_{pq})_{r\times t},}
且对应块同型，则
\fml{A^{\mathrm{T}}=(A_{pq}^{\mathrm{T}})_{t\times r},\qquad kA=(kA_{pq})_{r\times t},\qquad A\pm B=(A_{pq}\pm B_{pq})_{r\times t}.}
对乘法，设 $A=(A_{pq})_{s\times t}$，$B=(B_{pq})_{t\times r}$，且 $A$ 的\textbf{列}的分法与 $B$ 的\textbf{行}的分法一致，则
\fml{AB=(C_{pq})_{s\times r},\qquad C_{pq}=\sum_{l=1}^{t}A_{pl}B_{lq}=A_{p1}B_{1q}+A_{p2}B_{2q}+\cdots+A_{pt}B_{tq}.}
关键之处在于：分块乘法与普通矩阵乘法\textbf{运算规则完全一致}（"小块当成元素"），
但相容性条件从"内阶相等"升级为"\textbf{分法一致}"——左矩阵按列分、右矩阵按行分，切口必须对得上。

\textbf{把块当数的唯一代价就是检查分法相容}。凡是分块矩阵相乘的题，先画虚线验证切口，
再按普通乘法写出 $C_{pq}$；若切口不齐，答案必然错且往往错在"少了若干块"。
转置是把"块的位置转置"再"每块自己转置"，这两步都要做，缺一即错。
【反哺点】服务考纲「矩阵」中"掌握分块矩阵及其运算"。这条是分块题的第一道门槛：
\end{fdeep}

% ---------------- ⑫ 准对角矩阵 ----------------
\begin{fdeep}{准对角矩阵：行列式等于各块行列式之积、逆等于各块之逆}
形为
\fml{A=\begin{pmatrix}A_1&&O\\&\ddots&\\O&&A_l\end{pmatrix}}
的矩阵，其中 $A_i\ (i=1,2,\dots,l)$ 是 $n_i\times n_i$ 矩阵，通常称为\textbf{准对角矩阵}
（分块对角矩阵），特别地当 $A_i$ 都是 $1$ 阶时就是\textbf{对角矩阵}。

（樊 4.3.2）分块对角矩阵的两条基本公式：
\fml{|A|=|A_1||A_2|\cdots|A_l|;\qquad\text{若各}\ A_i\ \text{可逆，则}\ A\ \text{可逆，且}\
A^{-1}=\begin{pmatrix}A_1^{-1}&&O\\&\ddots&\\O&&A_l^{-1}\end{pmatrix}.}
（北大 §5）对于两个有相同分块的准对角矩阵，如果它们相应的分块是同阶的，则显然有
\fml{AB=\begin{pmatrix}A_1B_1&&O\\&\ddots&\\O&&A_lB_l\end{pmatrix},\qquad
A+B=\begin{pmatrix}A_1+B_1&&O\\&\ddots&\\O&&A_l+B_l\end{pmatrix},}
它们还是准对角矩阵；再结合上面的逆公式，\textbf{准对角矩阵类在加法、乘法、可逆求逆下封闭}。

（a）行列式可以各块各算再相乘——遇到数值上明显分块对角的行列式，先分块再降阶；
（b）求逆可以各块求逆——高阶求逆题若矩阵是准对角的，直接对小块用二阶求逆公式；
（c）分块运算保结构——证明题中"该类矩阵对运算封闭"一句话即可，不必逐元素验证。
注意逆公式与乘法公式的顺序：因为各块彼此不干扰，这里的顺序问题自动消失（这正是准对角的特殊性）。
【反哺点】服务考纲「矩阵」中"掌握分块矩阵及其运算"。这三条公式是分块题里"\textbf{能拆就拆}"的底气：
\end{fdeep}

% ---------------- ⑬ 分块上三角求逆与 Schur 补 ----------------
\begin{fdeep}{分块上（下）三角的逆：Schur 补从哪来}
（北大 §5）设 $D=\begin{pmatrix}A&O\\C&B\end{pmatrix}$，其中 $A,B$ 分别是 $k$ 阶和 $r$ 阶的\textbf{可逆}矩阵，
$C$ 是 $r\times k$ 矩阵。首先 $\begin{vmatrix}A&O\\C&B\end{vmatrix}=|A|\,|B|$，所以 $D$ 可逆。
设 $D^{-1}=\begin{pmatrix}X_{11}&X_{12}\\X_{21}&X_{22}\end{pmatrix}$，由分块乘法并比较等式两边，得
\fml{\begin{cases}AX_{11}=E_k,\quad AX_{12}=O,\\ CX_{11}+BX_{21}=O,\quad CX_{12}+BX_{22}=E_r,\end{cases}}
由此 $X_{11}=A^{-1}$，$X_{12}=O$，$X_{22}=B^{-1}$，代回第三式得 $BX_{21}=-CX_{11}=-CA^{-1}$，
$X_{21}=-B^{-1}CA^{-1}$。因此
\fml{D^{-1}=\begin{pmatrix}A^{-1}&O\\-B^{-1}CA^{-1}&B^{-1}\end{pmatrix}.}

这段的价值有二：（a）给出\textbf{分块求逆的标准流程}——设未知分块、乘开、比较四组等式、
按"能直接读出的先读"的顺序回代，四步适用于任何分块形状；
（b）揭示了"$B^{-1}CA^{-1}$"这种\textbf{三明治结构}的来源——它来自 $BX_{21}=-CA^{-1}$ 这一步，
是分块求逆最容易记错正负号的位置。另注意 $|D|=|A||B|$ 与 $C$ 无关：
块三角的行列式只由对角块决定，这是 $\begin{vmatrix}A&O\\C&B\end{vmatrix}$ 型题的一行解法。
【反哺点】服务考纲「矩阵」中"掌握分块矩阵及其运算"（分块矩阵求逆是高频题型）。
\end{fdeep}

% ---------------- ⑭ 分块初等变换 ----------------
\begin{fdeep}{分块初等变换：左乘一个分块矩阵就做一次「分块行变换」}
（北大 §7）将单位矩阵分块为 $\begin{pmatrix}E_m&O\\O&E_n\end{pmatrix}$，对它进行两行（列）互换、
某一行（列）左乘（右乘）一个矩阵 $P$、一行（列）加上另一行（列）的 $P$（矩阵）倍数，
就得到分块初等矩阵。和初等矩阵与初等变换的关系一样，用它们\textbf{左}乘任一个分块矩阵
$\begin{pmatrix}A&B\\C&D\end{pmatrix}$，只要分块乘法能够进行，结果就是对它作相应的变换：
\fml{\begin{pmatrix}O&E_n\\E_m&O\end{pmatrix}\begin{pmatrix}A&B\\C&D\end{pmatrix}=\begin{pmatrix}C&D\\A&B\end{pmatrix},}
\fml{\begin{pmatrix}P&O\\O&E_n\end{pmatrix}\begin{pmatrix}A&B\\C&D\end{pmatrix}=\begin{pmatrix}PA&PB\\C&D\end{pmatrix},\qquad
\begin{pmatrix}E_m&O\\P&E_n\end{pmatrix}\begin{pmatrix}A&B\\C&D\end{pmatrix}=\begin{pmatrix}A&B\\C+PA&D+PB\end{pmatrix}.}
同样，用它们右乘任一矩阵，进行分块乘法时也有相应的结果（即对列作相应的分块变换）。
【反哺点】服务考纲「矩阵」中"掌握分块矩阵及其运算"。这三张表是分块技巧的\textbf{工具箱}，
用法是先想"我要把哪个块变成零或变成想要的形状"，再查表选矩阵：要交换两块用
$\begin{pmatrix}O&E\\E&O\end{pmatrix}$；要给某块左乘 $P$ 用 $\begin{pmatrix}P&O\\O&E\end{pmatrix}$
（注意它同时作用于该块所在的两列）；要把左下块 $C$ 消掉就用 $\begin{pmatrix}E&O\\P&E\end{pmatrix}$，
$P$ 由 $C+PA=O$ 定出，即 $P=-CA^{-1}$（$A$ 可逆时）。
这与第 ⑬ 块的分块求逆是同一件事的"变换语言"版本，也是下一块分块消去法的全部工具。
\end{fdeep}

% ---------------- ⑮ 分块消去法 ----------------
\begin{fdeep}{分块消去法（Schur 补）：一行一列地造零块}
（北大 §7）在 ⑭ 的 (3) 式中适当选择 $P$，可使 $C+PA=O$。例如 $A$ 可逆时选 $P=-CA^{-1}$，则 $C+PA=O$，
于是 (3) 的右端成为
\fml{\begin{pmatrix}A&B\\O&D-CA^{-1}B\end{pmatrix},}
这种形状的矩阵在求行列式、逆矩阵和解决其他问题时是比较方便的。
（北大 §7 例 3 的思路）由此可以证明行列式的乘积公式 $|A||B|=|AB|$：设 $A,B$ 为 $n\times n$ 矩阵，作
\fml{\begin{pmatrix}E_n&A\\O&E_n\end{pmatrix}\begin{pmatrix}A&O\\-E_n&B\end{pmatrix}=\begin{pmatrix}O&AB\\-E_n&B\end{pmatrix}.}
左端的左因子可写成 $P_{11}P_{12}\cdots P_{nn}$ 之积，而这类 $P_{ij}$ 所对应的初等变换是
"某行加上另外一行的倍数"，\textbf{它不改变行列式的值}；右端经 $n$ 个两列对换与块三角公式即得 $|AB|=|A||B|$。

"\textbf{哪个块可逆，就用它把对面的块消掉}"：消左下块用左乘 $\begin{pmatrix}E&O\\-CA^{-1}&E\end{pmatrix}$，
消完剩下的对角块 $D-CA^{-1}B$（Schur 补）自动携带全部"耦合信息"。
这套手法一次通吃三类考题：\textbf{算分块行列式}（用 $|D|=|A|\,|D-CA^{-1}B|$ 降阶）、
\textbf{证可逆性}（Schur 补可逆 $\Rightarrow$ 原矩阵可逆）、\textbf{证秩的等式}
（消去不改变秩，把矩阵化为分块三角后秩一目了然）。
例 3 还顺带说明：$|AB|=|A||B|$ 可以用"\textbf{分块造零块 $+$ 列对换 $+$ 不改变行列式的行变换}"三步证出来，
比按排列展开更短，是"用分块初等变换证行列式等式"的样板。
【反哺点】服务考纲「矩阵」中"掌握分块矩阵及其运算"。分块消去法的口诀是
\end{fdeep}

% ---------------- ⑯ 秩的乘法不等式 ----------------
\begin{fdeep}{秩的乘法不等式 $\operatorname{r}(AB)\le\min\{\operatorname{r}(A),\operatorname{r}(B)\}$（定理 2）}
\textbf{定理 2}（北大 §3）：设 $A$ 是数域 $P$ 上 $n\times m$ 矩阵，$B$ 是 $P$ 上 $m\times s$ 矩阵，于是
\fml{\operatorname{r}(AB)\le\min\{\operatorname{r}(A),\operatorname{r}(B)\},}
即乘积的秩不超过各因子的秩。

证明思路（用分块乘法，一端一步）：设 $B_1,B_2,\dots,B_m$ 表示 $B$ 的行向量。
把 $B$ \textbf{按行分块}，按分块相乘就有
\fml{AB=\begin{pmatrix}a_{11}B_1+a_{12}B_2+\cdots+a_{1m}B_m\\ \cdots\cdots\cdots\cdots\\
a_{n1}B_1+a_{n2}B_2+\cdots+a_{nm}B_m\end{pmatrix}.}
这个式子表明，矩阵 $AB$ 的行向量组可以经矩阵 $B$ 的行向量组\textbf{线性表出}，因而前者的秩不可能超过后者的秩：
$\operatorname{r}(AB)\le\operatorname{r}(B)$。同样，把 $B$ 按列分块（用 $A$ 的列向量 $A_1,\dots,A_m$）
可知 $AB$ 的列向量组可经 $A$ 的列向量组线性表出，故 $\operatorname{r}(AB)\le\operatorname{r}(A)$。合起来即得结论。
用数学归纳法不难推广到多个因子的情形：$\operatorname{r}(A_1A_2\cdots A_k)\le\min\limits_{1\le i\le k}\operatorname{r}(A_i)$。
【反哺点】服务考纲「矩阵」与「矩阵的秩」（第 4 讲的前置）。这条不等式是不等式类题的\textbf{起点}：
（a）见到"$AB=O$"立刻想到 $\operatorname{r}(A)+\operatorname{r}(B)\le n$（因 $AB$ 的列都是 $Ax=0$ 的解，
$r(B)\le n-r(A)$），这是齐次方程组与秩联动的核心；
（b）见到"$AB=E$"则 $\operatorname{r}(A)\ge n$，逼出 "$A$ 行满秩、$B$ 列满秩"；
（c）证明思路本身要记住——\textbf{把秩的比较翻译成"一个向量组能否由另一个线性表出"}，
这是所有秩不等式证明的统一手法（第 4 讲的 Sylvester 不等式也走这条路）。
\end{fdeep}

% ---------------- ⑰ 可逆矩阵不改变秩 ----------------
\begin{fdeep}{可逆矩阵不改变秩：$\operatorname{r}(A)=\operatorname{r}(PA)=\operatorname{r}(AQ)$（定理 4）}
\textbf{定理 4}（北大 §3）：$A$ 是一个 $s\times n$ 矩阵，如果 $P$ 是 $s\times s$ 可逆矩阵，
$Q$ 是 $n\times n$ 可逆矩阵，那么
\fml{\operatorname{r}(A)=\operatorname{r}(PA)=\operatorname{r}(AQ).}
证明思路（夹逼）：令 $B=PA$，由定理 2 得
\fml{\operatorname{r}(B)\le\operatorname{r}(A);}
但是由 $A=P^{-1}B$，又有 $\operatorname{r}(A)\le\operatorname{r}(B)$。所以
$\operatorname{r}(A)=\operatorname{r}(B)=\operatorname{r}(PA)$。另一个等式同样地证明。

这是"初等变换可以求秩"的\textbf{理论依据}（初等矩阵可逆，故左乘右乘都不改秩），
同时给出秩的\textbf{三大不变量}结论：左乘或右乘可逆矩阵后秩不变，于是
$\operatorname{r}(PAQ)=\operatorname{r}(A)$，即\textbf{等价矩阵秩相同}（与 ⑧、⑩ 的等价判据互为一体）。
实战用法：（a）题目里出现"乘了一堆可逆矩阵"时，直接把它们当空气，只盯原矩阵的秩；
（b）证明"$P$ 可逆时 $\operatorname{r}(PA)=\operatorname{r}(A)$"这类结论，
标准三步是"用 $\le$ 得一半 $+$ 用逆矩阵把另一半倒过来 $+$ 夹逼"，与 ④ 的等价刻画网同源。
【反哺点】服务考纲「矩阵」中"掌握用初等变换求矩阵的秩和逆矩阵的方法"与"了解矩阵等价"。
\end{fdeep}


% ===========================================================================
%  【主控裁决 · 第 3 章深化料取舍】（依据 AGENTS.md §2 决定二）
%
%  与主控自撰块 deep_ch03_authored.tex 的重复，删除下列块：
%   · 块1「矩阵乘法不可交换是常态」→ 自撰块「为什么"矩阵不能像数一样运算"」已覆盖，
%       且自撰版额外给出 C(A) 的子空间结构与 dim C(A)≥n 的判据。
%   · 块3「方阵乘积的行列式与可逆性：退化性可遗传可传染」→ 自撰块已涉及；
%       其 |AB|=|A||B| 部分属第 1 章内容。
%   · 块4「可逆的等价刻画网」→ 自撰块「可逆的等价刻画：把七条串成一张网」更完整
%       （含 (A|E)→(E|A^{-1}) 的由来与"只验一边"的应用）。
%   · 块6「初等矩阵的定义与三类形状（定义 10）」→ 属正文级基础定义，非深化；
%       块7「左乘做行变换、右乘做列变换」已含其实质。
%
%  迁移到第 4 章：
%   · 块16「秩的乘法不等式 r(AB)≤min{r(A),r(B)}」与块17「可逆矩阵不改变秩」
%       → 属第 4 章「矩阵的秩」主题（OUTLINE §5 的 4.3/4.5），在第 3 章会造成主题越界。
%         块17 暂留本章（服务"矩阵的等价"），块16 移至第 4 章。
%
%  保留 12 块：2,5,7,8,9,10,11,12,13,14,15,17
% ===========================================================================
